\subsection{Measure supports and singular launch laws}
\label{subsec:measure-support}
For a nonnegative locally integrable function $f$, the convention
used here is
\[
 \operatorname{supp}(f(x)\dd x)
 =\left\{x:\int_{B(x,r)}f(y)\dd y>0
                         \text{ for every }r>0\right\}.
\]
This set depends only on the almost-everywhere class of $f$.
It need not be the pointwise positivity set.  We write
$\check\mu=(-\operatorname{id})_\#\mu$.

\begin{lemma}[Positive convolution supports]
\label{lem:positive-convolution-support}
Let $\alpha,\beta$ be nonzero, nonnegative, compactly supported
finite Borel measures on $\R^2$.  Then
\begin{equation}\label{eq:positive-convolution-support}
 \operatorname{supp}(\alpha*\beta)
            =\operatorname{supp}\alpha+\operatorname{supp}\beta.
\end{equation}
If bounded measurable sets $E,E'$ differ by a planar null set,
then $\one_E*\beta=\one_{E'}*\beta$ almost everywhere,
including when $\beta$ is singular.
\end{lemma}
\begin{proof}
Put $U=\operatorname{supp}\alpha$ and
$V=\operatorname{supp}\beta$.  The convolution is carried by the
compact set $U+V$.  Conversely, if $x=p+q$ with $p\in U$ and
$q\in V$, then, for every $r>0$,
\[
 (\alpha*\beta)(B(x,r))
 \ge \alpha(B(p,r/3))\beta(B(q,r/3))>0.
\]
This proves \eqref{eq:positive-convolution-support}.
Tonelli's theorem identifies the convolution
$(\one_E(x)\dd x)*\beta$ with the measure having density
$\one_E*\beta$.  It also gives
\[
 \int_{\R^2}|\one_E*\beta-\one_{E'}*\beta|\dd x
 \le \beta(\R^2)\,\area(E\mathbin\triangle E')=0.
\]
Neither argument requires a density for either measure.
\end{proof}

\begin{proposition}[Observable components under an arbitrary compact law]
\label{prop:singular-support-components}
Under \eqref{eq:two-field-gap}, let
$v_C(x)=\int\one_C(x+z)\dd\mu(z)$, and let
$f_{\pm,C}(x)=\int B^{C}_{\pm a}(x+z)\dd\mu(z)$,
where $B^{C}$ is the collision bit for the single obstacle $C$.
Use coordinates $y=ue+we^\perp$ and write the sections of $C$ as
$[\ell_C(w),u_C(w)]$, with $w\in I_C=[w_-,w_+]$.
Define the closed arcs and strips by
\begin{align*}
 \Gamma_{+,C}&=\{\ell_C(w)e+we^\perp:w\in I_C\},&
 E_{+,C}&=\Gamma_{+,C}+[-t,0]e,\\
 \Gamma_{-,C}&=\{u_C(w)e+we^\perp:w\in I_C\},&
 E_{-,C}&=\Gamma_{-,C}+[0,t]e.
\end{align*}
Then
\begin{equation}\label{eq:singular-individual-supports}
 \operatorname{supp}(v_C(x)\dd x)=P_C:=C+(-A),\qquad
 \operatorname{supp}(f_{\pm,C}(x)\dd x)
                         =K_{\pm,C}:=E_{\pm,C}+(-A).
\end{equation}
The connected components of the occupation support are precisely
the $P_C$; those of the two collision supports are, respectively,
the $K_{+,C}$ and the $K_{-,C}$.  Moreover,
\begin{align}
 \dist(P_C,P_{C'})&\ge d-\Delta,\label{eq:singular-occupation-gap}\\
 \dist(K_{\pm,C},K_{\pm,C'})&\ge d-t-\Delta,
       \qquad C\ne C',\label{eq:singular-collision-gap}\\
 K_{\pm,C}\cap P_C&\ne\varnothing,\qquad
 \dist(K_{\pm,C},P_{C'})\ge d-t-\Delta
       \quad(C\ne C').\label{eq:singular-support-matching}
\end{align}
Thus intersection with the occupation support determines the
correspondence between the two collision component families.
These statements apply to point, segment, atomic and singular
continuous launch laws with the stated convex support.
\end{proposition}
\begin{proof}
The section functions $\ell_C$ and $u_C$ are convex and concave,
respectively, on $I_C$.  They are continuous in its interior.
Strict convexity makes each extreme section a singleton.  Any
convergent subsequence of section endpoints as $w\to w_\pm$
lies in that singleton, by compactness of $C$.  Both section
functions are consequently continuous on $I_C$.

For the positive command, the single-obstacle collision-start set
is exactly
\[
 D_{+,C}=\{(\ell_C(w)-s)e+we^\perp:
                         w\in I_C,\ 0<s\le t\}.
\]
The endpoint $s=0$ is a solid start and is excluded; $s=t$ is
a collision at the terminal endpoint and is included.  At
$w=w_\pm$ the same formula includes a tangential encounter.
The open set
\[
 G_{+,C}=\{(\ell_C(w)-s)e+we^\perp:
                    w_-<w<w_+,\ 0<s<t\}
\]
has closure $E_{+,C}$.  Indeed, the parametrization is a
homeomorphism from $I_C\times[0,t]$ to $E_{+,C}$, and its
restriction to the open rectangle has open image.  The sets
$D_{+,C}$ and $G_{+,C}$ differ only on two graphs and two
transverse endpoint segments, all of which are planar null sets
by one-dimensional slicing.  The same description, with
$u_C(w)+s$ in place of $\ell_C(w)-s$, applies to the negative
command.  In particular,
\[
 \operatorname{supp}(\one_{G_{\pm,C}}(x)\dd x)=E_{\pm,C}.
\]
Every neighborhood of a point of $E_{\pm,C}$ meets its open
interior in positive area, including at a tangent endpoint.

Since $C$ is the closure of its interior,
$\operatorname{supp}(\one_C(x)\dd x)=C$.  Apply
Lemma~\ref{lem:positive-convolution-support} to these three
Lebesgue measures and to $\check\mu$, whose support is $-A$.
The null-set assertion in that lemma passes from $G_{\pm,C}$
to the physical collision bits after convolution and proves
\eqref{eq:singular-individual-supports}.  This passage is an
almost-everywhere assertion; it does not discard the prescribed
pointwise boundary convention for the physical means.

For completeness, the occupation is positive at every point of
$\operatorname{int}P_C$.  The convex-set identity needed here is
\begin{equation}\label{eq:singular-interior-sum}
             \operatorname{int}(C+(-A))
                         =\operatorname{int}C+(-A).
\end{equation}
To prove the nontrivial inclusion, fix $c_0\in\operatorname{int}C$
and $a_0\in A$.  If $x\in\operatorname{int}(C-A)$, then for
small $\lambda>0$ the point
$(x-\lambda(c_0-a_0))/(1-\lambda)$ belongs to $C-A$.
Write it as $c-a$.  The identity
\[
 x=((1-\lambda)c+\lambda c_0)
                  -((1-\lambda)a+\lambda a_0)
\]
places $x$ in $\operatorname{int}C-A$; the reverse inclusion
follows because the latter set is a union of open translates.
If $x=c-a$ with $c\in\operatorname{int}C$ and $a\in A$, a
sufficiently small neighborhood of $a$ is mapped by $z\mapsto x+z$
into $C$.  That neighborhood has positive $\mu$-mass, so
$v_C(x)>0$.  Outside $P_C$ the occupation is zero.  Its positive
set therefore contains $\operatorname{int}P_C$ and is contained
in $P_C$; it is path connected, since a segment joining any of
its points to an interior point stays in the interior except
possibly at its first endpoint.  Its closure is $P_C$.

Each $E_{\pm,C}$ is path connected by its rectangle
parametrization.  Convexity of $A$ makes $K_{\pm,C}$ path
connected as well.  The estimates
\eqref{eq:singular-occupation-gap}--\eqref{eq:singular-support-matching}
follow by subtracting two representations of their points:
obstacle points have distance at least $d$, the difference of
two same-sign longitudinal shifts has length at most $t$, and
the difference of two footprint shifts has length at most
$\Delta$.  The nonempty intersection follows from
\[
 \Gamma_{\pm,C}+(-A)\subset K_{\pm,C}\cap P_C.
\]

All expanded families remain locally finite: a member meeting a
fixed compact set comes from an obstacle meeting its enlargement
by the fixed bounded set $A+[-t,t]e$.  Since $t<d$, a segment
cannot meet two different obstacles or start in one obstacle and
reach another.  Thus $v=\sum_Cv_C$ and
$F_\pm=\sum_C f_{\pm,C}$, with the physical boundary convention.
Their measure supports are the locally finite unions in
\eqref{eq:singular-individual-supports}.  The positive gaps and
connectedness just proved identify their connected components
and establish the matching assertion.
\end{proof}

For example, for any nonempty compact convex $A$, choose a sequence
$(a_j)_{j\ge1}$ dense in $A$ and put
$\mu=\sum_{j\ge1}2^{-j}\delta_{a_j}$.
Its topological support is $A$, and the proposition applies to this
purely atomic probability. This includes convex supports of dimension
zero, one or two and illustrates why the measure-support convention
is used throughout.

\begin{corollary}[Almost-everywhere data and the period group]
\label{cor:singular-support-ae-periods}
The almost-everywhere classes of $F_+,F_-$ determine the same
occupation and collision support components, with the same
matching, as their physical representatives.  If
\[
 \operatorname{Per}_{\mathrm{ae}}(F_+,F_-)
 =\{w:F_\pm(\cdot+w)=F_\pm\text{ almost everywhere}\},
\]
then
\begin{equation}\label{eq:singular-support-periods}
 \operatorname{Per}_{\mathrm{ae}}(F_+,F_-)
 =\operatorname{Per}_{\mathrm{point}}(F_+,F_-)
 =\operatorname{Per}(\mathcal O).
\end{equation}
\end{corollary}
\begin{proof}
The finite prefix formula in Lemma~\ref{lem:two-field-prefix}
uses only finitely many fixed translations.  A change of the
fields on null sets changes its output only on a finite union
of translated null sets.  It therefore recovers the
almost-everywhere occupation class and hence its measure support.
The two original measure supports and their incidence relation
are likewise unchanged.  Proposition~\ref{prop:singular-support-components}
identifies the components without any choice of representatives.

If $w$ is an almost-everywhere common field period, the same
prefix formula gives $v(\cdot+w)=v$ almost everywhere.
Translation by $w$ permutes the components of its measure
support, so $P_C+w=P_{C'}$ for some $C'$.  Taking support
functions and canceling the common $h_{-A}$ gives
$h_C(u)+w\cdot u=h_{C'}(u)$ for every unit $u$.
Hence $C+w=C'$ and $\mathcal O+w=\mathcal O$.
Conversely, a period of $\mathcal O$ preserves every physical
collision bit pointwise, and integration against the same law
preserves every physical mean pointwise.  This proves
\eqref{eq:singular-support-periods}.
\end{proof}
